\chapter{लुइस संरचनाएँ, अनुनाद और VSEPR}\label{ch:b1:lewis-resonance-vsepr}

ओज़ोन अणु, \ce{O3}, एक मुड़ी हुई श्रृंखला में तीन ऑक्सीजन परमाणु है। इसे
पुस्तक 1 के नियमों से बनाइए तो एक आबंध द्वि-आबंध आता है और दूसरा
एकल: एक आबंध छोटा होना चाहिए, दूसरा लंबा। फिर भी सूक्ष्मतरंग
स्पेक्ट्रमिकी ओज़ोन के दोनों \ce{O-O} आबंधों को ठीक बराबर पाती है,
\qty{127.8}{pm}, जो \ce{O2} के द्वि-आबंध (\qty{120.8}{pm}) और
हाइड्रोजन परॉक्साइड के एकल आबंध (\qty{147.5}{pm}) के बीच है। कोई एक चित्र
सही नहीं है; दो चित्र मिलकर सही हैं। यह अध्याय लुइस मॉडल को
सटीक बनाता है --- \omterm{def:b1:lewis-resonance-vsepr:formal-charge}{औपचारिक आवेश}, अपूर्ण या अतिरिक्त अष्टक,
कई संरचनाओं के बीच अनुनाद --- और फिर उससे अणुओं की आकृतियों और उनसे
निकलने वाले \omterm{def:b1:lewis-resonance-vsepr:dipole}{द्विध्रुव आघूर्णों} का पूर्वानुमान लगाता है।
% ledger: geo:O3.rOO, re:O2, geo:H2O2.rOO

\begin{recall}
पुस्तक 1 (कक्षा 10) ने \omterm{def:b1:lewis-resonance-vsepr:lewis}{सहसंयोजी आबंध} को इलेक्ट्रॉनों के एक साझा युग्म के रूप में
बताया था, आबंधी और \omterm{def:b1:lewis-resonance-vsepr:lewis}{एकाकी युग्मों} के साथ \omterm{def:b1:lewis-resonance-vsepr:lewis}{लुइस संरचनाएँ} बनाई थीं, और
चार आकृतियों (रैखिक, कोणीय, त्रिकोणीय, चतुष्फलकीय) का पूर्वानुमान लगाया था।
\omterm{def:b1:periodicity:electronegativity}{विद्युत-ऋणात्मकता} और उसके अंतरों को \cref{ch:b1:periodicity} में मात्रात्मक
बनाया गया; किसी परमाणु के \omterm{def:b1:quantum-numbers:configuration}{संयोजकता इलेक्ट्रॉन}
उसके विन्यास से पढ़े जाते हैं (\cref{def:b1:quantum-numbers:configuration})।
\end{recall}

\section{लुइस संरचनाएँ और औपचारिक आवेश}

\begin{definition}[सहसंयोजी आबंध, लुइस संरचना]\label{def:b1:lewis-resonance-vsepr:lewis}
\emph{सहसंयोजी आबंध}\index{सहसंयोजी आबंध} दो परमाणुओं द्वारा साझा किया गया
इलेक्ट्रॉनों का एक युग्म है, जिसे \emph{आबंधी युग्म}\index{आबंधी युग्म} कहते हैं; दो या तीन
युग्म द्वि-आबंध या त्रि-आबंध बनाते हैं। \omterm{def:b1:quantum-numbers:configuration}{संयोजकता इलेक्ट्रॉनों} का वह युग्म जो
किसी एक ही परमाणु का हो, \emph{एकाकी युग्म}\index{एकाकी युग्म} है।
\emph{लुइस संरचना}\index{लुइस संरचना} किसी अणु या आयन के हर \omterm{def:b1:quantum-numbers:configuration}{संयोजकता
इलेक्ट्रॉन} को दिखाती है, आबंधों (परमाणुओं के बीच रेखाएँ) और एकाकी
युग्मों (किसी परमाणु पर रेखाएँ या बिंदुओं के जोड़े) के रूप में। \emph{अष्टक
नियम}\index{अष्टक नियम} के अनुसार \omterm{def:b1:periodicity:table}{आवर्त} 2 के परमाणु (C, N, O, F) किसी
स्थायी संरचना में चार युग्मों, यानी आठ इलेक्ट्रॉनों, से घिरे होते हैं;
\emph{द्विक नियम}\index{द्विक नियम} के अनुसार हाइड्रोजन एक युग्म से।
\end{definition}

\begin{definition}[औपचारिक आवेश]\label{def:b1:lewis-resonance-vsepr:formal-charge}
\omterm{def:b1:lewis-resonance-vsepr:lewis}{लुइस संरचना} में किसी परमाणु का \emph{औपचारिक आवेश}\index{औपचारिक आवेश}
है
\[
  q_F = N_v - N_{\text{एकाकी}} - \tfrac12 N_{\text{आबंध}},
\]
जहाँ $N_v$ मुक्त परमाणु के \omterm{def:b1:quantum-numbers:configuration}{संयोजकता इलेक्ट्रॉनों} की संख्या है,
$N_{\text{एकाकी}}$ उसके \omterm{def:b1:lewis-resonance-vsepr:lewis}{एकाकी युग्मों} में इलेक्ट्रॉनों की संख्या और
$N_{\text{आबंध}}$ उसके द्वारा बनाए गए आबंधों में इलेक्ट्रॉनों की संख्या। शून्य न होने पर
इसे परमाणु के पास $\oplus$ या $\ominus$ के रूप में लिखा जाता है।
\end{definition}

\begin{proposition}[औपचारिक आवेशों का योग आवेश के बराबर है]\label{prop:b1:lewis-resonance-vsepr:formal-sum}
किसी \omterm{def:b1:lewis-resonance-vsepr:lewis}{लुइस संरचना} के परमाणुओं के \omterm{def:b1:lewis-resonance-vsepr:formal-charge}{औपचारिक आवेशों} का योग
अणु या आयन के कुल आवेश के बराबर होता है।
\end{proposition}

\begin{proof}
परमाणुओं पर $q_F$ का योग करने से $\sum N_v - (\text{एकाकी
युग्मों के इलेक्ट्रॉन} + \text{आबंधों के इलेक्ट्रॉन})$ मिलता है, क्योंकि हर आबंध के इलेक्ट्रॉन
उसके दोनों परमाणुओं पर आधे-आधे गिने जाते हैं। कोष्ठक संरचना में बनाए गए
\omterm{def:b1:quantum-numbers:configuration}{संयोजकता इलेक्ट्रॉनों} की कुल संख्या है, जो $\sum N_v$ में से
स्पीशीज़ का आवेश $z$ घटाने पर मिलती है। अतः $\sum q_F = z$।
\end{proof}

\begin{method}[लुइस संरचना बनाना]\label{met:b1:lewis-resonance-vsepr:lewis}
किसी अणु या आयन की \omterm{def:b1:lewis-resonance-vsepr:lewis}{लुइस संरचना} बनाने के लिए:
\begin{enumerate}
  \item \omterm{def:b1:quantum-numbers:configuration}{संयोजकता इलेक्ट्रॉन} $N = \sum N_v - z$ और युग्म $N/2$ गिनिए;
  \item ढाँचा बनाइए: सबसे कम विद्युत-ऋणात्मक परमाणु (H नहीं) प्रायः
        केंद्रीय होता है; हाइड्रोजन और फ़्लुओरीन सदा सिरे पर होते हैं;
  \item सिरे वाले परमाणुओं के अष्टक \omterm{def:b1:lewis-resonance-vsepr:lewis}{एकाकी युग्मों} से पूरे कीजिए, फिर
        बचे हुए युग्म केंद्रीय परमाणु पर रखिए;
  \item यदि केंद्रीय परमाणु का अष्टक अधूरा है, तो उसके पड़ोसियों के \omterm{def:b1:lewis-resonance-vsepr:lewis}{एकाकी युग्मों} को
        बहु-आबंधों में बदलिए;
  \item \omterm{def:b1:lewis-resonance-vsepr:formal-charge}{औपचारिक आवेश} परिकलित कीजिए; संभव संरचनाओं में उसे वरीयता दीजिए
        जिसमें सबसे कम \omterm{def:b1:lewis-resonance-vsepr:formal-charge}{औपचारिक आवेश} हों, और ऋणात्मक आवेश
        सबसे अधिक विद्युत-ऋणात्मक परमाणुओं पर हों।
\end{enumerate}
\end{method}

\begin{example}[नाइट्रिक अम्ल]\label{ex:b1:lewis-resonance-vsepr:hno3}
\ce{HNO3}: $N = 1 + 5 + 3 \times 6 = 24$ इलेक्ट्रॉन, 12 युग्म। ढाँचा:
नाइट्रोजन केंद्र में, तीन ऑक्सीजनों से आबंधित, जिनमें से एक पर
हाइड्रोजन है। एकल आबंधों से अष्टक पूरे करने पर नाइट्रोजन के पास छह
इलेक्ट्रॉन बचते हैं; ऑक्सीजन का एक \omterm{def:b1:lewis-resonance-vsepr:lewis}{एकाकी युग्म} \ce{N=O} आबंध बन जाता है। \omterm{def:b1:lewis-resonance-vsepr:formal-charge}{औपचारिक आवेश}:
नाइट्रोजन $5 - 0 - 4 = +1$; हाइड्रोजन-रहित एकल-आबंधित ऑक्सीजन
$6 - 6 - 1 = -1$; बाक़ी 0। कुल 0 है, जैसा
\cref{prop:b1:lewis-resonance-vsepr:formal-sum} माँगती है।
\end{example}

\begin{omfigure}
\setchemfig{atom sep=2.4em}
\begin{tikzpicture}[font=\small]
  \node at (0,0) {\large\chemfig{H-\charge{90=\|,270=\|}{O}-\charge{90:2pt=$\scriptstyle\oplus$}{N}(=[:45]\charge{0=\|,90=\|}{O})-[:-45]\charge{0=\|,240=\|,300=\|,45:3pt=$\scriptstyle\ominus$}{O}}};
  \node[below] at (0,-1.5) {नाइट्रिक अम्ल, \ce{HNO3}};
  \node at (5.2,0) {\large\chemfig{\charge{180=\|}{N}~\charge{0=\|}{N}}};
  \node[below] at (5.2,-1.5) {डाइनाइट्रोजन, \ce{N2}};
  \node at (9.6,0) {\large\chemfig{\charge{135=\|,225=\|,90:2pt=$\scriptstyle\ominus$}{N}=\charge{90:2pt=$\scriptstyle\oplus$}{N}=\charge{45=\|,315=\|}{O}}};
  \node[below] at (9.6,-1.5) {डाइनाइट्रोजन मोनॉक्साइड, \ce{N2O}};
\end{tikzpicture}

{\small \omterm{def:b1:lewis-resonance-vsepr:lewis}{एकाकी युग्मों} और \omterm{def:b1:lewis-resonance-vsepr:formal-charge}{औपचारिक आवेशों} सहित \omterm{def:b1:lewis-resonance-vsepr:lewis}{लुइस संरचनाएँ}।}
\end{omfigure}

\section{अष्टक से आगे}

\begin{definition}[इलेक्ट्रॉन-न्यून और अतिसंयोजी अणु]\label{def:b1:lewis-resonance-vsepr:hypervalent}
जिस अणु में किसी परमाणु के पास आठ से कम \omterm{def:b1:quantum-numbers:configuration}{संयोजकता इलेक्ट्रॉन} हों, वह
\emph{इलेक्ट्रॉन-न्यून}\index{इलेक्ट्रॉन-न्यून} है: \ce{BF3} में बोरॉन
के पास छह हैं। जिस अणु में \omterm{def:b1:periodicity:table}{आवर्त} 3 या उसके आगे का कोई परमाणु चार से अधिक युग्मों से
घिरा हो, वह \emph{अतिसंयोजी}\index{अतिसंयोजी} है:
\ce{PCl5} में फ़ॉस्फ़ोरस पाँच आबंध बनाता है, \ce{SF6} में सल्फ़र छह। विषम संख्या
में इलेक्ट्रॉनों वाले अणु, जैसे \ce{NO} और \ce{NO2}, हर परमाणु पर \omterm{def:b1:lewis-resonance-vsepr:lewis}{अष्टक नियम}
पूरा नहीं कर सकते: एक इलेक्ट्रॉन अयुग्मित रहता है।
\end{definition}

\begin{remark}[आवर्त 3 क्यों और आवर्त 2 क्यों नहीं]\label{rem:b1:lewis-resonance-vsepr:hypervalent}
नाइट्रोजन \ce{NF3} बनाता है पर कभी \ce{NF5} नहीं, जबकि फ़ॉस्फ़ोरस
\ce{PF3} और \ce{PF5} दोनों बनाता है। \omterm{def:b1:periodicity:table}{आवर्त} 2 का परमाणु चार से अधिक पड़ोसियों
से आबंधित होने के लिए बहुत छोटा है; \omterm{def:b1:periodicity:table}{आवर्त} 3 का परमाणु पाँच या छह को थामने लायक बड़ा है।
\ce{SO4^{2-}} या \ce{H2SO4} के चित्र सल्फ़र के द्वि-आबंधों के साथ (S के चारों ओर 12
इलेक्ट्रॉन) या एकल आबंधों और \omterm{def:b1:lewis-resonance-vsepr:formal-charge}{औपचारिक आवेशों} के साथ (एक अष्टक), दोनों तरह
मिलते हैं; दूसरा आवेश-वितरण को बेहतर दर्शाता है, और
दोनों सही ज्यामिति देते हैं।
\end{remark}

\begin{definition}[लुइस अम्ल, लुइस क्षारक, उपसहसंयोजी आबंध]\label{def:b1:lewis-resonance-vsepr:lewis-acid}
\emph{लुइस अम्ल}\index{लुइस अम्ल} वह स्पीशीज़ है जो अपने संयोजकता कोश के
किसी रिक्त स्थान में इलेक्ट्रॉनों का एक युग्म ग्रहण कर सकती है (\ce{BF3},
\ce{AlCl3}, \ce{H+}, धातु धनायन); \emph{लुइस क्षारक}\index{लुइस क्षारक}
वह स्पीशीज़ है जिसके पास साझा करने योग्य \omterm{def:b1:lewis-resonance-vsepr:lewis}{एकाकी युग्म} हो (\ce{NH3}, \ce{H2O},
\ce{F-})। जब क्षारक वह युग्म अम्ल को देता है, तो बनने वाला आबंध
\emph{उपसहसंयोजी आबंध}\index{उपसहसंयोजी आबंध} होता है: ऐसा \omterm{def:b1:lewis-resonance-vsepr:lewis}{सहसंयोजी आबंध} जिसके दोनों
इलेक्ट्रॉन एक ही साथी से आए हों।
\end{definition}

\begin{example}[अमोनिया और बोरॉन ट्राइफ़्लुओराइड]\label{ex:b1:lewis-resonance-vsepr:adduct}
\ce{NH3 + BF3 -> H3NBF3}। योगोत्पाद में बोरॉन अपना अष्टक पूरा करता है;
नाइट्रोजन, जो अब अपना \omterm{def:b1:lewis-resonance-vsepr:lewis}{एकाकी युग्म} साझा करता है, पर \omterm{def:b1:lewis-resonance-vsepr:formal-charge}{औपचारिक आवेश}
$+1$ है और बोरॉन पर $-1$। एक बार बन जाने पर \ce{N-B} आबंध एक साधारण
\omterm{def:b1:lewis-resonance-vsepr:lewis}{सहसंयोजी आबंध} है।
\end{example}

\begin{omfigure}
\setchemfig{atom sep=2.1em}
\begin{tikzpicture}[font=\small]
  \node at (0,0) {\large\chemfig{H-\charge{90=\|}{N}(-[:240]H)(-[:300]H)}};
  \node at (1.6,0) {$+$};
  \node at (3.1,0) {\large\chemfig{B(-[:90]F)(-[:210]F)(-[:330]F)}};
  \draw[-{Stealth}, thick] (4.3,0) -- (5.3,0);
  \node at (7.4,0) {\large\chemfig{H-\charge{45:1pt=$\scriptstyle\oplus$}{N}(-[:90]H)(-[:270]H)-\charge{45:1pt=$\scriptstyle\ominus$}{B}(-[:90]F)(-[:270]F)-F}};
\end{tikzpicture}

{\small एक \omterm{def:b1:lewis-resonance-vsepr:lewis-acid}{उपसहसंयोजी आबंध}: अमोनिया (\omterm{def:b1:lewis-resonance-vsepr:lewis-acid}{लुइस क्षारक}) का \omterm{def:b1:lewis-resonance-vsepr:lewis}{एकाकी युग्म}
बोरॉन ट्राइफ़्लुओराइड (\omterm{def:b1:lewis-resonance-vsepr:lewis-acid}{लुइस अम्ल}) का रिक्त स्थान भरता है। फ़्लुओरीन के \omterm{def:b1:lewis-resonance-vsepr:lewis}{एकाकी युग्म}
नहीं दिखाए गए हैं।}
\end{omfigure}

\section{अनुनाद}

\begin{definition}[अनुनादी संरचनाएँ, अनुनाद संकर]\label{def:b1:lewis-resonance-vsepr:resonance}
जब किसी अणु के इलेक्ट्रॉन कई ऐसी \omterm{def:b1:lewis-resonance-vsepr:lewis}{लुइस संरचनाओं} में रखे जा सकें जो केवल
$\pi$ आबंधों और \omterm{def:b1:lewis-resonance-vsepr:lewis}{एकाकी युग्मों} की स्थिति में भिन्न हों --- नाभिक अपनी जगह
रहते हुए --- तो हर एक \emph{अनुनादी
संरचना}\index{अनुनादी संरचना} (या मेसोमेरी रूप) है, जो
दूसरों से दो सिरों वाले तीर $\leftrightarrow$ द्वारा जुड़ी होती है। वास्तविक अणु
इनमें से कोई नहीं है: वह \emph{अनुनाद संकर}\index{अनुनाद संकर} है,
एक अकेली संरचना जिसका इलेक्ट्रॉन-वितरण उनके वितरणों का भारित
औसत है।
\end{definition}

\begin{remark}[एक तीर, साम्य नहीं]\label{rem:b1:lewis-resonance-vsepr:arrow}
तीर $\leftrightarrow$ किसी अभिक्रिया का वर्णन नहीं करता: ओज़ोन एक संरचना से
दूसरी में बदलता नहीं रहता। संकर एक ही अणु है, जिसका वर्णन कई
चित्रों से इसलिए किया जाता है कि एक चित्र पर्याप्त नहीं। \omterm{def:b1:lewis-resonance-vsepr:resonance}{अनुनादी संरचनाओं} के
बीच कभी $\rightleftharpoons$ का प्रयोग मत कीजिए।
\end{remark}

\begin{method}[अनुनादी संरचनाएँ लिखना और तौलना]\label{met:b1:lewis-resonance-vsepr:resonance}
किसी स्पीशीज़ की \omterm{def:b1:lewis-resonance-vsepr:resonance}{अनुनादी संरचनाएँ} खोजने के लिए:
\begin{enumerate}
  \item एक \omterm{def:b1:lewis-resonance-vsepr:lewis}{लुइस संरचना} से शुरू कीजिए; ऐसा \omterm{def:b1:lewis-resonance-vsepr:lewis}{एकाकी युग्म} या
        $\pi$ आबंध खोजिए जो किसी $\pi$ आबंध, रिक्त स्थान या धनात्मक
        आवेश के बगल में हो;
  \item युग्मों को वक्र तीरों से खिसकाइए (\omterm{def:b1:lewis-resonance-vsepr:lewis}{एकाकी युग्म} $\pi$ आबंध बनता है,
        $\pi$ आबंध \omterm{def:b1:lewis-resonance-vsepr:lewis}{एकाकी युग्म} बनता है), कभी किसी नाभिक को खिसकाए बिना और कभी
        \omterm{def:b1:periodicity:table}{आवर्त} 2 के परमाणु का अष्टक पार किए बिना;
  \item \omterm{def:b1:lewis-resonance-vsepr:formal-charge}{औपचारिक आवेश} फिर से परिकलित कीजिए;
  \item संरचनाओं को तौलिए: सबसे महत्वपूर्ण संरचनाएँ \omterm{def:b1:lewis-resonance-vsepr:lewis}{अष्टक
        नियम} का पालन करती हैं, उनमें सबसे कम \omterm{def:b1:lewis-resonance-vsepr:formal-charge}{औपचारिक आवेश} होते हैं, और ऋणात्मक आवेश
        सबसे अधिक विद्युत-ऋणात्मक परमाणुओं पर होते हैं; तुल्य संरचनाओं का भार
        बराबर होता है।
\end{enumerate}
\end{method}

\begin{example}[ओज़ोन, नाइट्रेट, बेंज़ीन]\label{ex:b1:lewis-resonance-vsepr:examples}
ओज़ोन की दो तुल्य संरचनाएँ हैं: हर \ce{O-O} आबंध एक में द्वि-आबंध है
और दूसरी में एकल, इसलिए दोनों की आबंध कोटि $1.5$ है और लंबाई
एक ही, \qty{127.8}{pm}, जो \ce{O=O} और \ce{O-O} के बीच है। नाइट्रेट आयन
\ce{NO3-} की तीन तुल्य संरचनाएँ हैं; हर \ce{N-O} आबंध की कोटि
$4/3$ है। बेंज़ीन, \ce{C6H6}, की दो केकुले संरचनाएँ हैं; उसके छहों \ce{C-C}
आबंध बराबर हैं, \qty{139.7}{pm}, जो एथेन के एकल आबंध
(\qty{153.6}{pm}) और एथीन के द्वि-आबंध (\qty{133.9}{pm}) के बीच है।
% ledger: geo:O3.rOO, geo:C6H6.rCC, geo:C2H6.rCC, geo:C2H4.rCC
\end{example}

\begin{omfigure}
\vspace*{10pt}
\large
\setchemfig{atom sep=2.4em}
\schemestart
\chemfig{\charge{135=\|,225=\|}{O}=[:-30]\charge{270=\|,90:2pt=$\scriptstyle\oplus$}{O}-[:30]\charge{90=\|,0=\|,270=\|,45:5pt=$\scriptstyle\ominus$}{O}}
\arrow{<->}
\chemfig{\charge{90=\|,180=\|,270=\|,135:5pt=$\scriptstyle\ominus$}{O}-[:-30]\charge{270=\|,90:2pt=$\scriptstyle\oplus$}{O}=[:30]\charge{45=\|,315=\|}{O}}
\schemestop

\medskip
\schemestart
\chemfig{*6(=-=-=-)}
\arrow{<->}
\chemfig{*6(-=-=-=)}
\schemestop
\hspace{3em}
\chemfig{\charge{270:2pt=$\scriptstyle\oplus$}{N}(=[:90]\charge{45=\|,135=\|}{O})(-[:210]\charge{150=\|,240=\|,300=\|,90:2pt=$\scriptstyle\ominus$}{O})-[:330]\charge{30=\|,300=\|,240=\|,90:2pt=$\scriptstyle\ominus$}{O}}
\normalsize
\par\vspace{14pt}
{\small ओज़ोन में अनुनाद (दाईं ओर वाले ऑक्सीजन का एक \omterm{def:b1:lewis-resonance-vsepr:lewis}{एकाकी युग्म}
$\pi$ आबंध बन जाता है जबकि बाईं ओर का $\pi$ आबंध \omterm{def:b1:lewis-resonance-vsepr:lewis}{एकाकी
युग्म} बन जाता है) और बेंज़ीन में; तथा नाइट्रेट आयन की तीन तुल्य संरचनाओं में से
एक, जिसका आवेश तीनों ऑक्सीजनों में बँटा है।}
\end{omfigure}

\begin{definition}[$\sigma$ आबंध, $\pi$ आबंध]\label{def:b1:lewis-resonance-vsepr:sigma-pi}
एकल आबंध एक \emph{$\sigma$ आबंध}\index{सिग्मा आबंध@$\sigma$ आबंध} होता है:
इसका इलेक्ट्रॉन युग्म दोनों नाभिकों को जोड़ने वाली अक्ष के अनुदिश होता है, जो दो
कक्षकों के आमने-सामने अतिव्यापन से बनता है। बहु-आबंध के दूसरे और तीसरे आबंध
\emph{$\pi$ आबंध}\index{पाई आबंध@$\pi$ आबंध} होते हैं: उनके युग्म
अक्ष के दोनों ओर होते हैं, जो अक्ष के लंबवत दो p कक्षकों के पार्श्व
अतिव्यापन से बनते हैं। द्वि-आबंध एक $\sigma$ और एक $\pi$ आबंध है;
त्रि-आबंध एक $\sigma$ और दो $\pi$।
\end{definition}

\begin{omfigure}
\begin{tikzpicture}[font=\small]
  % sigma: two lobes overlapping head-on
  \fill[omDef!35, opacity=0.8] (0,0) ellipse (0.9 and 0.42);
  \fill[omProp!35, opacity=0.8] (1.1,0) ellipse (0.9 and 0.42);
  \fill (0.1,0) circle (1.6pt); \fill (1.0,0) circle (1.6pt);
  \draw[dashed, black!50] (-1.2,0) -- (2.3,0);
  \node[below] at (0.55,-1.05) {$\sigma$: आमने-सामने, अक्ष पर};
  % pi: two p orbitals side by side
  \begin{scope}[xshift=5.2cm]
    \foreach \x/\c in {0/omDef, 1.2/omProp} {
      \fill[\c!35, opacity=0.85] (\x,0.45) ellipse (0.32 and 0.45);
      \fill[\c!15, opacity=0.85] (\x,-0.45) ellipse (0.32 and 0.45);
      \draw[\c!70!black] (\x,0.45) ellipse (0.32 and 0.45) (\x,-0.45) ellipse (0.32 and 0.45);
      \fill (\x,0) circle (1.6pt);
    }
    \draw[dashed, black!50] (-0.7,0) -- (1.9,0);
    \draw[omThm, thick, densely dotted] (0.25,0.75) -- (0.95,0.75) (0.25,-0.75) -- (0.95,-0.75);
    \node[below] at (0.6,-1.05) {$\pi$: पार्श्व, ऊपर और नीचे};
  \end{scope}
\end{tikzpicture}

{\small अतिव्यापन के दो प्रकार। $\sigma$ आबंध अपनी अक्ष के चारों ओर घूर्णन की
अनुमति देता है; $\pi$ आबंध नहीं देता, इसीलिए द्वि-आबंध दृढ़ होता है
(\cref{ch:b1:stereochemistry-in-depth})।}
\end{omfigure}

\section{VSEPR मॉडल}

\begin{definition}[VSEPR मॉडल, त्रिविम संख्या]\label{def:b1:lewis-resonance-vsepr:vsepr}
\emph{VSEPR मॉडल}\index{VSEPR मॉडल} (संयोजकता-कोश
इलेक्ट्रॉन-युग्म प्रतिकर्षण) में किसी केंद्रीय परमाणु A के चारों ओर के युग्म एक-दूसरे को
प्रतिकर्षित करते हैं और वह व्यवस्था अपनाते हैं जिसमें वे एक-दूसरे से सबसे दूर रहें।
अणु को \ce{AX_nE_m} लिखा जाता है, जहाँ $n$, A से आबंधित परमाणुओं की
संख्या है (बहु-आबंध एक बार गिना जाता है) और $m$, A के \omterm{def:b1:lewis-resonance-vsepr:lewis}{एकाकी युग्मों} की
संख्या। \emph{त्रिविम संख्या}\index{त्रिविम संख्या} $n + m$ युग्मों की
व्यवस्था तय करती है: 2 रैखिक, 3 त्रिकोणीय समतलीय, 4 चतुष्फलकीय, 5
त्रिकोणीय द्विपिरामिडी, 6 अष्टफलकीय। अणु की आकृति केवल उसके
परमाणुओं की व्यवस्था है।
\end{definition}

\begin{proposition}[VSEPR ज्यामितियाँ]\label{prop:b1:lewis-resonance-vsepr:geometries}
\omterm{def:b1:lewis-resonance-vsepr:vsepr}{VSEPR मॉडल} नीचे की सारणी की आकृतियों और आदर्श कोणों का पूर्वानुमान लगाता है;
मापे गए कोण इनके पास हैं।
\begin{center}
\small
\begin{tabular}{llllc}
\hline
प्रकार & युग्मों की व्यवस्था & आकृति & उदाहरण & मापा गया कोण \\
\hline
\ce{AX2} & रैखिक & रैखिक & \ce{CO2} & $180^\circ$ \\
\ce{AX3} & त्रिकोणीय समतलीय & त्रिकोणीय समतलीय & \ce{BF3} & $120^\circ$ \\
\ce{AX2E} & त्रिकोणीय समतलीय & कोणीय & \ce{SO2} & $119.5^\circ$ \\
\ce{AX4} & चतुष्फलकीय & चतुष्फलकीय & \ce{CH4} & $109.5^\circ$ \\
\ce{AX3E} & चतुष्फलकीय & त्रिकोणीय पिरामिडी & \ce{NH3} & $106.7^\circ$ \\
\ce{AX2E2} & चतुष्फलकीय & कोणीय & \ce{H2O} & $104.5^\circ$ \\
\ce{AX5} & त्रिकोणीय द्विपिरामिडी & त्रिकोणीय द्विपिरामिडी & \ce{PCl5} & $90^\circ$, $120^\circ$ \\
\ce{AX4E} & त्रिकोणीय द्विपिरामिडी & ढेंकली & \ce{SF4} & $101.6^\circ$, $173.1^\circ$ \\
\ce{AX3E2} & त्रिकोणीय द्विपिरामिडी & T-आकार & \ce{ClF3} & $87.5^\circ$ \\
\ce{AX2E3} & त्रिकोणीय द्विपिरामिडी & रैखिक & \ce{XeF2} & $180^\circ$ \\
\ce{AX6} & अष्टफलकीय & अष्टफलकीय & \ce{SF6} & $90^\circ$ \\
\ce{AX5E} & अष्टफलकीय & वर्ग पिरामिडी & \ce{BrF5} & --- \\
\ce{AX4E2} & अष्टफलकीय & वर्ग समतलीय & \ce{XeF4} & --- \\
\hline
\end{tabular}
\end{center}
% ledger: geo:CO2.aOCO, geo:BF3.aFBF, geo:SO2.aOSO, geo:CH4.aHCH,
% geo:NH3.aHNH, geo:H2O.aHOH, geo:SF6.aFSF, geo:SF4.aFSF2, geo:SF4.aFSF,
% geo:ClF3.aFClF, geo:XeF2.aFXeF
\end{proposition}
\admitted

\begin{proposition}[एकाकी युग्म कोणों को छोटा करते हैं]\label{prop:b1:lewis-resonance-vsepr:lone-pair-angles}
\omterm{def:b1:lewis-resonance-vsepr:lewis}{एकाकी युग्म}, जिसे केवल एक नाभिक थामता है, \omterm{def:b1:lewis-resonance-vsepr:lewis}{आबंधी युग्म} से अधिक फैलता है
और अधिक प्रबलता से प्रतिकर्षित करता है: प्रतिकर्षण इस क्रम में घटते हैं: \omterm{def:b1:lewis-resonance-vsepr:lewis}{एकाकी
युग्म}/\omterm{def:b1:lewis-resonance-vsepr:lewis}{एकाकी युग्म} $>$ \omterm{def:b1:lewis-resonance-vsepr:lewis}{एकाकी युग्म}/\omterm{def:b1:lewis-resonance-vsepr:lewis}{आबंधी युग्म} $>$ \omterm{def:b1:lewis-resonance-vsepr:lewis}{आबंधी युग्म}/\omterm{def:b1:lewis-resonance-vsepr:lewis}{आबंधी युग्म}।
अतः \omterm{def:b1:lewis-resonance-vsepr:lewis}{एकाकी युग्म} जुड़ने के साथ आबंधों के बीच के कोण छोटे होते जाते हैं
(\ce{CH4} में $109.5^\circ$, \ce{NH3} में $106.7^\circ$,
\ce{H2O} में $104.5^\circ$), और त्रिकोणीय द्विपिरामिड में \omterm{def:b1:lewis-resonance-vsepr:lewis}{एकाकी युग्म} विषुवतीय
स्थान लेते हैं, जहाँ $90^\circ$ पर उनके केवल दो पड़ोसी होते हैं।
% ledger: geo:CH4.aHCH, geo:NH3.aHNH, geo:H2O.aHOH
\end{proposition}
\admitted

\begin{omfigure}
\setchemfig{atom sep=2.6em}
\renewcommand{\arraystretch}{1.2}
\begin{tabular}{ccccc}
\chemfig{O=C=O} &
\chemfig{B(-[:90]F)(-[:210]F)(-[:330]F)} &
\chemfig{C(-[:90]H)(-[:200]H)(<[:280]H)(<:[:335]H)} &
\chemfig{P(-[:90]Cl)(-[:270]Cl)(-[:180]Cl)(<[:320]Cl)(<:[:30]Cl)} &
\chemfig{S(-[:90]F)(-[:270]F)(-[:180]F)(-[:0]F)(<[:225]F)(<:[:45]F)} \\
\small \ce{AX2}, \ce{CO2} & \small \ce{AX3}, \ce{BF3} & \small \ce{AX4}, \ce{CH4} &
\small \ce{AX5}, \ce{PCl5} & \small \ce{AX6}, \ce{SF6} \\[16pt]
\chemfig{\charge{90=\:}{S}(=[:210]O)(=[:330]O)} &
\chemfig{\charge{90=\:}{N}(-[:200]H)(<[:280]H)(<:[:335]H)} &
\chemfig{\charge{60=\:,120=\:}{O}(-[:215]H)(-[:325]H)} &
\chemfig{\charge{0=\:}{S}(-[:90]F)(-[:270]F)(<[:210]F)(<:[:150]F)} &
\chemfig{Xe(-[:0]F)(-[:90]F)(-[:180]F)(-[:270]F)} \\
\small \ce{AX2E}, \ce{SO2} & \small \ce{AX3E}, \ce{NH3} & \small \ce{AX2E2}, \ce{H2O} &
\small \ce{AX4E}, \ce{SF4} & \small \ce{AX4E2}, \ce{XeF4} \\
\end{tabular}

\medskip
{\small VSEPR द्वारा अनुमानित आकृतियाँ। वेज पाठक की ओर इंगित करता है,
धारीदार वेज दूर; सादे आबंध पृष्ठ के तल में हैं। केंद्रीय परमाणु के \omterm{def:b1:lewis-resonance-vsepr:lewis}{एकाकी
युग्म} बिंदुओं के जोड़ों के रूप में दिखाए गए हैं; \ce{XeF4} सामने से
बनाया गया है, उसके दो \omterm{def:b1:lewis-resonance-vsepr:lewis}{एकाकी युग्म} पाठक की ओर और उससे दूर इंगित करते हैं।}
\end{omfigure}

\begin{method}[आकृति का पूर्वानुमान]\label{met:b1:lewis-resonance-vsepr:vsepr}
किसी केंद्रीय परमाणु के चारों ओर की आकृति का पूर्वानुमान लगाने के लिए:
\begin{enumerate}
  \item \omterm{def:b1:lewis-resonance-vsepr:lewis}{लुइस संरचना} बनाइए (\cref{met:b1:lewis-resonance-vsepr:lewis});
  \item केंद्रीय परमाणु से आबंधित परमाणु ($n$) और उसके \omterm{def:b1:lewis-resonance-vsepr:lewis}{एकाकी
        युग्म} ($m$) गिनिए; \ce{AX_nE_m} लिखिए;
  \item $n + m$ से व्यवस्था पढ़िए और $n$ परमाणुओं की
        स्थितियों से आकृति;
  \item आदर्श कोणों को सुधारिए: \omterm{def:b1:lewis-resonance-vsepr:lewis}{एकाकी युग्म} और बहु-आबंध एकल आबंधों से अधिक
        जगह लेते हैं।
\end{enumerate}
\end{method}

\begin{definition}[संकरण]\label{def:b1:lewis-resonance-vsepr:hybridisation}
किसी परमाणु का \emph{संकरण}\index{संकरण} उसके आबंधों का ऐसे कक्षकों
द्वारा वर्णन है जो VSEPR द्वारा अनुमानित दिशाओं में इंगित करते हैं:
\omterm{def:b1:lewis-resonance-vsepr:vsepr}{त्रिविम संख्या} 4 वाले परमाणु को sp$^3$ कहा जाता है
($109.5^\circ$ पर चार तुल्य कक्षक), \omterm{def:b1:lewis-resonance-vsepr:vsepr}{त्रिविम संख्या} 3 वाले को sp$^2$
(एक तल में $120^\circ$ पर तीन कक्षक, और उस तल के लंबवत एक p कक्षक
बचा हुआ), \omterm{def:b1:lewis-resonance-vsepr:vsepr}{त्रिविम संख्या} 2 वाले को sp ($180^\circ$ पर दो कक्षक,
दो p कक्षक बचे हुए)। बचे हुए p कक्षक $\pi$ आबंध बनाते हैं।
\end{definition}

\begin{example}[कार्बनिक रसायन के कार्बन]\label{ex:b1:lewis-resonance-vsepr:carbon}
एथेन में हर कार्बन sp$^3$, चतुष्फलकीय है; एथीन में sp$^2$, छहों
परमाणु एक तल में, कोण $120^\circ$ के पास (मापा गया \ce{H-C-H}
$117.6^\circ$); एथाइन में sp, चारों परमाणु एक रेखा पर। $\pi$ आबंध जुड़ने के साथ आबंध
लंबाइयाँ 153.6 से 133.9 और फिर \qty{120.3}{pm} तक घटती हैं।
% ledger: geo:C2H4.aHCH, geo:C2H6.rCC, geo:C2H4.rCC, geo:C2H2.rCC
\end{example}

\begin{remark}[एक वर्णन, कारण नहीं]\label{rem:b1:lewis-resonance-vsepr:hybrid-limit}
\omterm{def:b1:lewis-resonance-vsepr:hybridisation}{संकरण} पहले से ज्ञात ज्यामिति का वर्णन करता है; वह उसकी व्याख्या नहीं
करता, और इलेक्ट्रॉनों की ऊर्जाओं के लिए विफल रहता है, जिनका हिसाब आण्विक कक्षक
सिद्धांत, द्वितीय वर्ष के खंड में, देता है। फिर भी यह कार्बनिक रसायन की
रोज़मर्रा की भाषा है: ``sp$^2$ कार्बन'' का अर्थ है एक समतलीय
कार्बन जिसके पास $\pi$ आबंध के लिए एक p कक्षक ख़ाली है।
\end{remark}

\section{ध्रुवीय अणु: द्विध्रुव आघूर्ण}

\begin{definition}[आबंध द्विध्रुव, द्विध्रुव आघूर्ण]\label{def:b1:lewis-resonance-vsepr:dipole}
भिन्न \omterm{def:b1:periodicity:electronegativity}{विद्युत-ऋणात्मकताओं} वाले परमाणुओं के बीच आबंध में इलेक्ट्रॉन
अधिक विद्युत-ऋणात्मक परमाणु की ओर झुकते हैं, जिस पर आंशिक आवेश
$-\delta e$ होता है जबकि उसके साथी पर $+\delta e$ ($0 < \delta < 1$)।
\emph{आबंध द्विध्रुव}\index{आबंध द्विध्रुव} सदिश $\vec\mu = \delta e\,
\vec d$ है, जिसका मान $\delta e d$ है और जो ऋणात्मक से धनात्मक
आंशिक आवेश की ओर इंगित करता है, जहाँ $\vec d$ उनके बीच का सदिश है। किसी अणु का
\emph{द्विध्रुव आघूर्ण}\index{द्विध्रुव आघूर्ण} उसके आबंध द्विध्रुवों (और उसके
\omterm{def:b1:lewis-resonance-vsepr:lewis}{एकाकी युग्मों} के योगदानों) का सदिश योग है; इसे
कूलॉम मीटर या डिबाई में व्यक्त किया जाता है, $\qty{1}{\debye} =
\qty{3.336e-30}{C.m}$। शून्येतर द्विध्रुव आघूर्ण वाला अणु
\emph{ध्रुवीय अणु}\index{ध्रुवीय अणु} है।
% ledger: const:c (1 D = 1e-21/c C.m)
\end{definition}

\begin{remark}[दिशा की परिपाटी]\label{rem:b1:lewis-resonance-vsepr:convention}
भौतिकी द्विध्रुव सदिश को ऋणात्मक से धनात्मक आवेश की ओर बनाती है,
और यहाँ यही परिपाटी अपनाई गई है। रसायन की पुरानी पुस्तकें कटी हुई पूँछ वाला एक तीर
बनाती हैं जो ऋणात्मक सिरे की ओर इंगित करता है; केवल
दिशा भिन्न है।
\end{remark}

\begin{proposition}[सममित अणु का द्विध्रुव]\label{prop:b1:lewis-resonance-vsepr:dipole-sum}
A पर \omterm{def:b1:lewis-resonance-vsepr:lewis}{एकाकी युग्मों} से रहित अणु \ce{AX_n} (रैखिक \ce{AX2}, त्रिकोणीय
\ce{AX3}, चतुष्फलकीय \ce{AX4}, \dots) का \omterm{def:b1:lewis-resonance-vsepr:dipole}{द्विध्रुव आघूर्ण} शून्य होता है, उसके
आबंधों की ध्रुवता चाहे जो हो। कोण $\theta$ वाले कोणीय \ce{AX2} अणु का,
जिसके आबंधों का द्विध्रुव $\mu_b$ है, \omterm{def:b1:lewis-resonance-vsepr:dipole}{द्विध्रुव आघूर्ण} है
\[
  \mu = 2\mu_b \cos\frac{\theta}{2} .
\]
\end{proposition}

\begin{proof}
सममित आकृतियों में आबंध सदिशों के मान बराबर होते हैं और उनकी
दिशाएँ A के चारों ओर सममित रूप से व्यवस्थित होती हैं, इसलिए उनका योग शून्य होता है
(रैखिक \ce{AX2} के लिए वे विपरीत हैं; $120^\circ$ वाले \ce{AX3} के लिए
समान मान के तीन सदिशों का योग शून्य है; \ce{AX4} के लिए योग
चतुष्फलक के घूर्णनों के अंतर्गत अपरिवर्तित रहता है, इसलिए शून्य है)। कोणीय
अणु के लिए हर आबंध सदिश समद्विभाजक से कोण $\theta/2$ बनाता है;
समद्विभाजक के लंबवत घटक निरस्त हो जाते हैं और उसके अनुदिश घटक
जुड़ जाते हैं: $2\mu_b\cos(\theta/2)$।
\end{proof}

\begin{omfigure}
\begin{tikzpicture}[font=\small, >={Stealth}]
  % water
  \node[aO] (o) at (0,0) {};
  \node[aH] (h1) at (-0.95,-0.73) {};
  \node[aH] (h2) at (0.95,-0.73) {};
  \begin{scope}[on background layer]
    \draw[bond] (o.center) -- (h1.center); \draw[bond] (o.center) -- (h2.center);
  \end{scope}
  \draw[->, omProp, thick] (-0.25,0.15) -- (-0.95,-0.38);
  \draw[->, omProp, thick] (0.25,0.15) -- (0.95,-0.38);
  \draw[->, omThm, very thick] (0,-0.25) -- (0,-1.45);
  \node[omThm, right] at (0.05,-1.25) {$\vec\mu$};
  \node at (0,0.55) {$\delta^-$};
  \node at (-1.25,-1.05) {$\delta^+$}; \node at (1.25,-1.05) {$\delta^+$};
  \node[below] at (0,-1.7) {\ce{H2O}: $\mu = \qty{1.86}{\debye}$};
  % carbon dioxide
  \begin{scope}[xshift=5cm]
    \node[aO] (a) at (-1.1,0) {}; \node[aC] (c) at (0,0) {}; \node[aO] (b) at (1.1,0) {};
    \begin{scope}[on background layer]
      \foreach \d in {-1.6,1.6} {
        \draw[bond, line width=1.1pt] ([yshift=\d pt]a.center) -- ([yshift=\d pt]c.center);
        \draw[bond, line width=1.1pt] ([yshift=\d pt]c.center) -- ([yshift=\d pt]b.center);}
    \end{scope}
    \draw[->, omProp, thick] (-0.75,0.45) -- (-0.2,0.45);
    \draw[->, omProp, thick] (0.75,0.45) -- (0.2,0.45);
    \node[below] at (0,-1.7) {\ce{CO2}: आबंध द्विध्रुव निरस्त, $\mu = 0$};
  \end{scope}
\end{tikzpicture}

{\small \omterm{def:b1:lewis-resonance-vsepr:dipole}{आबंध द्विध्रुव} (नारंगी, आंशिक ऋणात्मक से आंशिक
धनात्मक आवेश की ओर) और उनका योग (लाल)। कोणीय पानी में वे जुड़ते हैं; रैखिक
कार्बन डाइऑक्साइड में निरस्त हो जाते हैं।}
% ledger: dip:H2O
\end{omfigure}

\begin{example}[ध्रुवीय और अध्रुवीय अणु]\label{ex:b1:lewis-resonance-vsepr:dipoles}
\ce{CO2}, \ce{BF3}, \ce{CH4} और \ce{CCl4} के \omterm{def:b1:lewis-resonance-vsepr:dipole}{द्विध्रुव आघूर्ण} शून्य हैं।
\ce{H2O} (\qty{1.86}{\debye}), \ce{NH3} (\qty{1.48}{\debye}) और
\ce{SO2} (\qty{1.63}{\debye}) ध्रुवीय हैं। \ce{CH3Cl}, \ce{CH2Cl2},
\ce{CHCl3} के क्रम में आघूर्ण घटता है, 1.87, 1.62, \qty{1.04}{\debye}, और
\ce{CCl4} में लुप्त हो जाता है: \ce{C-Cl} \omterm{def:b1:lewis-resonance-vsepr:dipole}{आबंध द्विध्रुव} उत्तरोत्तर अधिक निरस्त होते हैं।
हाइड्रोजन हैलाइडों में यह \omterm{def:b1:periodicity:electronegativity}{विद्युत-ऋणात्मकता} के अंतर का अनुसरण करता है: HF 1.83,
HCl 1.09, HBr 0.83, HI \qty{0.45}{\debye}।
% ledger: dip:H2O, dip:NH3, dip:SO2, dip:CH3Cl, dip:CH2Cl2, dip:CHCl3,
% dip:CCl4, dip:HF, dip:HCl, dip:HBr, dip:HI
\end{example}

\section{अभ्यास}

\begin{exercise}[$\star$]\label{exo:b1:lewis-resonance-vsepr:1}
\ce{H2O}, \ce{NH3}, \ce{CO2}, \ce{HCN},
\ce{CH2O} (मेथेनैल) और \ce{NH4+} की \omterm{def:b1:lewis-resonance-vsepr:lewis}{लुइस संरचनाएँ} बनाइए, और \omterm{def:b1:lewis-resonance-vsepr:formal-charge}{औपचारिक आवेश} बताइए।
\end{exercise}

\begin{exercise}[$\star$]\label{exo:b1:lewis-resonance-vsepr:2}
इन स्पीशीज़ का प्रकार \ce{AX_nE_m}, आकृति और अनुमानित कोण बताइए:
\ce{H2S}, \ce{PCl3}, \ce{BF3}, \ce{SiH4} और \ce{CO3^{2-}} (हर सूत्र में केंद्रीय परमाणु
पहले)।
\end{exercise}

\begin{exercise}[$\star$]\label{exo:b1:lewis-resonance-vsepr:3}
इनमें से कौन-से अणु ध्रुवीय हैं: \ce{CCl4}, \ce{CH2Cl2}, \ce{BF3},
\ce{NH3}, \ce{CO2}, \ce{SO2}? आकृतियों से कारण बताइए।
\end{exercise}

\begin{exercise}[$\star$]\label{exo:b1:lewis-resonance-vsepr:4}
\ce{H+ + H2O -> H3O+} में और
\ce{Cu^{2+} + 4NH3 -> [Cu(NH3)4]^{2+}} में \omterm{def:b1:lewis-resonance-vsepr:lewis-acid}{लुइस अम्ल} और \omterm{def:b1:lewis-resonance-vsepr:lewis-acid}{लुइस क्षारक} पहचानिए। कौन-से आबंध उपसहसंयोजी हैं?
\end{exercise}

\begin{exercise}[$\star\star$]\label{exo:b1:lewis-resonance-vsepr:5}
एथेनोएट आयन \ce{CH3COO-} की दो \omterm{def:b1:lewis-resonance-vsepr:resonance}{अनुनादी संरचनाएँ} लिखिए। वे दोनों \ce{C-O}
आबंध लंबाइयों के बारे में क्या अनुमान देती हैं? मेथेनोइक
अम्ल \ce{HCOOH} से तुलना कीजिए, जिसके दो \ce{C-O} आबंध 120.2 और
\qty{134.3}{pm} के हैं।
% ledger: geo:H2CO2.rCO, geo:H2CO2.rCO2
\end{exercise}

\begin{exercise}[$\star\star$]\label{exo:b1:lewis-resonance-vsepr:6}
सल्फ़ेट आयन \ce{SO4^{2-}} की \omterm{def:b1:lewis-resonance-vsepr:lewis}{लुइस संरचना} चार
एकल \ce{S-O} आबंधों के साथ बनाइए, \omterm{def:b1:lewis-resonance-vsepr:formal-charge}{औपचारिक आवेश} परिकलित कीजिए, और
आकृति का अनुमान लगाइए। दो \ce{S=O} आबंधों वाली और सल्फ़र पर बिना
आवेश वाली कितनी \omterm{def:b1:lewis-resonance-vsepr:resonance}{अनुनादी संरचनाएँ} लिखी जा सकती हैं?
\end{exercise}

\begin{exercise}[$\star\star$]\label{exo:b1:lewis-resonance-vsepr:7}
VSEPR से \ce{SF4}, \ce{ClF3}, \ce{XeF2} और
\ce{XeF4} की आकृतियों का पूर्वानुमान लगाइए। समझाइए कि त्रिकोणीय द्विपिरामिड के \omterm{def:b1:lewis-resonance-vsepr:lewis}{एकाकी युग्म}
विषुवतीय तल में क्यों रहते हैं।
\end{exercise}

\begin{exercise}[$\star\star$]\label{exo:b1:lewis-resonance-vsepr:8}
कोण \ce{H-X-H}, \ce{H2O} में $104.5^\circ$ और
\ce{H2S} में $92.1^\circ$ हैं; \ce{NH3} में $106.7^\circ$ और \ce{PH3} में $93.3^\circ$। केंद्रीय
परमाणु के आकार और \omterm{def:b1:periodicity:electronegativity}{विद्युत-ऋणात्मकता} के आधार पर एक व्याख्या
प्रस्तावित कीजिए।
% ledger: geo:H2O.aHOH, geo:H2S.aHSH, geo:NH3.aHNH, geo:PH3.aHPH
\end{exercise}

\begin{exercise}[$\star\star$]\label{exo:b1:lewis-resonance-vsepr:9}
\ce{CH3-C#N} (एथेननाइट्राइल) में और \ce{CH2=CH-CH3} में हर कार्बन और
नाइट्रोजन का \omterm{def:b1:lewis-resonance-vsepr:hybridisation}{संकरण} बताइए, और हर अणु में
$\sigma$ और $\pi$ आबंधों की संख्या भी।
\end{exercise}

\begin{exercise}[$\star\star\star$]\label{exo:b1:lewis-resonance-vsepr:10}
\ce{NF3} का \omterm{def:b1:lewis-resonance-vsepr:dipole}{द्विध्रुव आघूर्ण} केवल \qty{0.24}{\debye} है, जबकि
\ce{NH3} का \qty{1.48}{\debye}, यद्यपि \ce{N-F} आबंध
\ce{N-H} आबंध से अधिक ध्रुवीय है। \omterm{def:b1:lewis-resonance-vsepr:dipole}{आबंध द्विध्रुवों} और नाइट्रोजन के
\omterm{def:b1:lewis-resonance-vsepr:lewis}{एकाकी युग्म} की सहायता से यह अंतर समझाइए।
% ledger: dip:NF3, dip:NH3
\end{exercise}

\begin{exercise}[$\star\star\star$]\label{exo:b1:lewis-resonance-vsepr:11}
\ce{SO2} अणु का कोण $119.5^\circ$ है और \omterm{def:b1:lewis-resonance-vsepr:dipole}{द्विध्रुव आघूर्ण}
\qty{1.63}{\debye}। एक \ce{S-O} आबंध का द्विध्रुव परिकलित कीजिए; आबंध
लंबाई \qty{143.2}{pm} लेकर हर ऑक्सीजन पर आंशिक आवेश $\delta$
निकालिए।
% ledger: geo:SO2.aOSO, geo:SO2.rSO, dip:SO2
\end{exercise}

\begin{exercise}[$\star\star\star$]\label{exo:b1:lewis-resonance-vsepr:12}
डाइनाइट्रोजन मोनॉक्साइड \ce{N2O}
(रैखिक, \ce{N-N-O}) की सबसे महत्वपूर्ण \omterm{def:b1:lewis-resonance-vsepr:lewis}{लुइस संरचनाएँ} बनाइए और \omterm{def:b1:lewis-resonance-vsepr:formal-charge}{औपचारिक आवेश} परिकलित कीजिए। मापी गई आबंध
लंबाइयाँ \ce{N-N} \qty{112.8}{pm} और \ce{N-O} \qty{118.4}{pm} हैं;
\ce{N2} में त्रि-आबंध \qty{109.8}{pm} का है। किस संरचना का भार
अधिक है?
% ledger: geo:N2O.rNN, geo:N2O.rNO, re:N2
\end{exercise}

\section{समस्या: तड़ित के अणु}

\begin{problem}[{सप्ताहांत समस्या --- तड़ित द्वारा वायु के गर्म होने पर बनने वाली
नाइट्रोजन और ऑक्सीजन की स्पीशीज़: लुइस संरचनाएँ, अनुनाद, आकृतियाँ, और
पानी के आंशिक आवेश}]
\label{pb:b1:lewis-resonance-vsepr:1}
तड़ित वायु को हज़ारों डिग्री तक गर्म कर देती है: \ce{N2} और \ce{O2}
से \ce{NO} बनता है, फिर \ce{NO2}, \ce{N2O}, ओज़ोन और अंत में वर्षा में नाइट्रिक अम्ल।
आँकड़े: $\qty{1}{\debye} = \qty{3.336e-30}{C.m}$, $e =
\qty{1.602e-19}{C}$। जहाँ आवश्यक है वहाँ मापे गए मान दिए गए हैं।

\medskip
\noindent\textbf{भाग I --- \omterm{def:b1:lewis-resonance-vsepr:lewis}{लुइस संरचनाएँ}।}
\begin{enumerate}
  \item \ce{N2}, \ce{NO}, \ce{NO2},
        \ce{N2O}, \ce{O3} और \ce{HNO3} के \omterm{def:b1:quantum-numbers:configuration}{संयोजकता इलेक्ट्रॉन} गिनिए।
  \item \ce{N2} की \omterm{def:b1:lewis-resonance-vsepr:lewis}{लुइस संरचना} बनाइए।
  \item \ce{NO} की एक \omterm{def:b1:lewis-resonance-vsepr:lewis}{लुइस संरचना} बनाइए। दोनों परमाणुओं पर \omterm{def:b1:lewis-resonance-vsepr:lewis}{अष्टक नियम}
        क्यों पूरा नहीं हो सकता?
  \item \ce{N2O} की दो \omterm{def:b1:lewis-resonance-vsepr:lewis}{लुइस संरचनाएँ} बनाइए, एक \ce{N=N=O} के साथ और एक
        \ce{N#N-O} के साथ, और \omterm{def:b1:lewis-resonance-vsepr:formal-charge}{औपचारिक आवेश} बताइए।
  \item ओज़ोन की एक \omterm{def:b1:lewis-resonance-vsepr:lewis}{लुइस संरचना} बनाइए और \omterm{def:b1:lewis-resonance-vsepr:formal-charge}{औपचारिक आवेश} बताइए।
  \item \ce{HNO3} की एक \omterm{def:b1:lewis-resonance-vsepr:lewis}{लुइस संरचना} बनाइए और \omterm{def:b1:lewis-resonance-vsepr:formal-charge}{औपचारिक आवेश} बताइए।
  \item छह स्पीशीज़ में से किनमें \omterm{def:b1:quantum-numbers:unpaired}{अयुग्मित इलेक्ट्रॉन} है?
\end{enumerate}

\medskip
\noindent\textbf{भाग II --- अनुनाद और आबंध लंबाइयाँ।}
\begin{enumerate}[resume]
  \item ओज़ोन की दो \omterm{def:b1:lewis-resonance-vsepr:resonance}{अनुनादी संरचनाएँ} लिखिए। वे कौन-सी आबंध कोटि
        बताती हैं?
  \item \ce{O-O} आबंध ओज़ोन में \qty{127.8}{pm} का है,
        \ce{O2} में \qty{120.8}{pm} और \ce{H2O2} में \qty{147.5}{pm} का। क्या
        पूर्वानुमान की पुष्टि होती है?
  \item \ce{NO3-} की तीन \omterm{def:b1:lewis-resonance-vsepr:resonance}{अनुनादी संरचनाएँ} लिखिए और हर \ce{N-O} आबंध की
        कोटि बताइए।
  \item \ce{HNO3} में तीन \ce{N-O} आबंध 140.6, 121.1 और
        \qty{119.9}{pm} के हैं। इन्हें निर्दिष्ट कीजिए और अनुनाद से समझाइए।
  \item समझाइए कि कोण \ce{O-N-O}, जो \ce{NO2} में $134^\circ$ है,
        $120^\circ$ से बड़ा क्यों है।
  \item प्रश्न 4 की दोनों संरचनाओं में से कौन-सी ऋणात्मक
        \omterm{def:b1:lewis-resonance-vsepr:formal-charge}{औपचारिक आवेश} को अधिक विद्युत-ऋणात्मक परमाणु पर रखती है?
\end{enumerate}

\medskip
\noindent\textbf{भाग III --- आकृतियाँ और ध्रुवता।}
\begin{enumerate}[resume]
  \item प्रकार \ce{AX_nE_m} बताइए: \ce{O3},
        \ce{N2O}, \ce{NO3-} के केंद्रीय परमाणु का और \ce{HNO3} के नाइट्रोजन का;
        साथ में उनकी आकृतियाँ भी।
  \item ओज़ोन का कोण $116.8^\circ$ है। समझाइए कि यह
        $120^\circ$ से कम क्यों है।
  \item ओज़ोन का \omterm{def:b1:lewis-resonance-vsepr:dipole}{द्विध्रुव आघूर्ण} \qty{0.53}{\debye} है, \ce{N2O} का
        \qty{0.16}{\debye}, \ce{N2} का कुछ नहीं। समझाइए।
  \item \ce{CO2} अध्रुवीय क्यों है जबकि \ce{SO2} ध्रुवीय है?
  \item $109.5^\circ$ की तुलना में \ce{NH3} के कोण का अनुमान लगाइए, फिर
        मापे गए $106.7^\circ$ से तुलना कीजिए।
  \item यही प्रश्न पानी के लिए ($104.5^\circ$)।
\end{enumerate}
% ledger: geo:O3.rOO, re:O2, geo:H2O2.rOO, geo:HNO3.rNO, geo:HNO3.rNO2,
% geo:HNO3.rNO3, geo:NO2.aONO, geo:O3.aOOO, dip:O3, dip:N2O

\medskip
\noindent\textbf{भाग IV --- पानी के आंशिक आवेश।}
पानी का \omterm{def:b1:lewis-resonance-vsepr:dipole}{द्विध्रुव आघूर्ण} \qty{1.857}{\debye} है, उसका कोण
$104.48^\circ$ और उसकी \ce{O-H} आबंध लंबाई \qty{95.8}{pm}।
% ledger: dip:H2O, geo:H2O.aHOH, geo:H2O.rOH
\begin{enumerate}[resume]
  \item पानी के \omterm{def:b1:lewis-resonance-vsepr:dipole}{द्विध्रुव आघूर्ण} को \omterm{def:b1:lewis-resonance-vsepr:dipole}{आबंध
        द्विध्रुव} $\mu_{OH}$ और कोण के फलन के रूप में व्यक्त कीजिए।
  \item $\mu_{OH}$ डिबाई में परिकलित कीजिए।
  \item इसे कूलॉम मीटर में बदलिए।
  \item यदि \omterm{def:b1:lewis-resonance-vsepr:dipole}{आबंध द्विध्रुव} दोनों नाभिकों पर दो पूर्ण आवेशों $\pm e$ से
        बना होता, तो वह कितना होता?
  \item \ce{O-H} आबंध का भिन्नात्मक आयनिक लक्षण $\delta$ निकालिए, यानी
        मापे गए \omterm{def:b1:lewis-resonance-vsepr:dipole}{आबंध द्विध्रुव} और पूर्ण
        आवेशों वाले द्विध्रुव का अनुपात।
\end{enumerate}
\end{problem}
